\section{Metodologia: protocolo da revisão sistemática}
\label{sec:metodologia}

A revisão segue as diretrizes de Kitchenham e Charters~\cite{kitchenham2007}, em
três fases (planeamento, condução e relato), e relata a seleção dos estudos com o
diagrama PRISMA 2020 (\emph{Preferred Reporting Items for Systematic Reviews and
Meta-Analyses})~\cite{page2021prisma}.

% ---------------------------------------------------------------------------
\subsection{Perguntas de investigação}
\label{sec:perguntas}

As perguntas (Tabela~\ref{tab:perguntas}) dividem-se numa pergunta geral, que
define \emph{o que} se investiga, em perguntas focadas, sobre \emph{como} os
problemas são tratados, e em perguntas estatísticas, que quantificam \emph{quanto}
é feito.

\begin{table}[ht]
  \centering
  \caption{Perguntas de investigação.}
  \label{tab:perguntas}
  \small
  \begin{tabularx}{\textwidth}{@{}L{1.1cm}Y@{}}
    \toprule
    \textbf{ID} & \textbf{Pergunta} \\
    \midrule
    QG  & Que mecanismos de segurança propõe a literatura para comunicações IoT
          baseadas em MQTT e que propriedades garantem? \\
    QF1 & Que ameaças e vulnerabilidades do MQTT são identificadas? \\
    QF2 & Que mecanismos são propostos e onde atuam (transporte, \emph{broker},
          mensagem fim-a-fim, deteção)? \\
    QF3 & Que propriedades garantem e assumem um \emph{broker} confiável? \\
    QF4 & Como são avaliadas as soluções? \\
    QF5 & Que limitações e desafios em aberto são reportados? \\
    QE1 & Quantas publicações existem por ano? \\
    QE2 & Qual a distribuição por tipo de publicação? \\
    QE3 & Quantos estudos abordam cada categoria e cada propriedade? \\
    \bottomrule
  \end{tabularx}
\end{table}

% ---------------------------------------------------------------------------
\subsection{Critérios PICOC}

\begin{table}[ht]
  \centering
  \caption{Âmbito da revisão segundo os critérios PICOC.}
  \label{tab:picoc}
  \small
  \begin{tabularx}{\textwidth}{@{}L{3cm}Y@{}}
    \toprule
    \textbf{Critério} & \textbf{Definição nesta revisão} \\
    \midrule
    \emph{Population}   & Comunicações IoT baseadas em MQTT ou MQTT-SN (MQTT para redes de sensores): clientes e
                          \emph{brokers}. \\
    \emph{Intervention} & Mecanismos de segurança: criptográficos, de autenticação e
                          autorização, de proteção do \emph{broker}, de deteção e de
                          privacidade. \\
    \emph{Comparison}   & MQTT com a segurança nativa ou apenas com TLS; mecanismos
                          propostos entre si. \\
    \emph{Outcome}      & Propriedades garantidas, ameaças mitigadas, custo e tipo de
                          avaliação. \\
    \emph{Context}      & Investigação com revisão por pares, de 2014 (normalização do
                          MQTT 3.1.1) a 2026. \\
    \bottomrule
  \end{tabularx}
\end{table}

% ---------------------------------------------------------------------------
\subsection{\emph{String} de pesquisa e bases de dados}
\label{sec:string}

A partir das perguntas identificaram-se grupos de ideias e os respetivos sinónimos
(Tabela~\ref{tab:grupos}); os termos de um grupo combinam-se com \texttt{OR} e os
grupos entre si com \texttt{AND}. A \emph{string} principal combina G1, G2 e G3 no
título, resumo e palavras-chave; a secundária acrescenta G4, para encontrar revisões
já publicadas. A \emph{string} principal inclui \texttt{"MQTT-SN"}, porque há
soluções para a variante para redes de sensores~\cite{singh2015secure,
sadio2019lightweight}, e \texttt{"intrusion"}, porque a deteção de intrusões é uma
linha própria~\cite{vaccari2020mqttset, alotaibi2024secure}.

\begin{table}[ht]
  \centering
  \caption{Grupos de ideias e sinónimos usados na \emph{string} de pesquisa.}
  \label{tab:grupos}
  \small
  \begin{tabularx}{\textwidth}{@{}L{2.8cm}Y@{}}
    \toprule
    \textbf{Grupo} & \textbf{Termos (combinados com OR)} \\
    \midrule
    G1 --- Protocolo   & \texttt{"MQTT"}, \texttt{"MQTT-SN"}, \texttt{"MQ Telemetry Transport"},
                         \texttt{"Message Queuing Telemetry Transport"},
                         \texttt{"Message Queue Telemetry Transport"} \\
    G2 --- Segurança   & \texttt{"security"}, \texttt{"secure"}, \texttt{"authentication"},
                         \texttt{"authorization"}, \texttt{"access control"},
                         \texttt{"encryption"}, \texttt{"confidentiality"},
                         \texttt{"integrity"}, \texttt{"privacy"}, \texttt{"attack*"},
                         \texttt{"vulnerabilit*"}, \texttt{"intrusion"} \\
    G3 --- Domínio     & \texttt{"Internet of Things"}, \texttt{"IoT"}, \texttt{"IIoT"},
                         \texttt{"publish/subscribe"}, \texttt{"publish-subscribe"},
                         \texttt{"machine-to-machine"}, \texttt{"M2M"} \\
    G4 --- Revisões \newline(só na secundária) & \texttt{"systematic literature"},
                         \texttt{"systematic review"}, \texttt{"literature review"},
                         \texttt{"survey"} \\
    \bottomrule
  \end{tabularx}
\end{table}

As pesquisas foram feitas no IEEE Xplore, no Scopus e no ScienceDirect, em
\dataPesquisa{}; a adaptação da \emph{string} a cada base consta do
Anexo~\ref{anx:strings}. O ScienceDirect, por não suportar o carácter \texttt{*} e
limitar os operadores booleanos, exigiu dividir a \emph{string} principal em duas
consultas, e a secundária, que excede esse limite, não pôde ser executada. O ACM
Digital Library e a Web of Science, previstos no protocolo, não foram pesquisados:
o primeiro exige subscrição para exportar resultados e a segunda não esteve
acessível; as atas ACM relevantes para o tema (por exemplo, ARES, SaT-CPS e ICPE)
estão indexadas no Scopus. Os registos foram exportados de cada base, com resumo, e
processados com um \emph{script} próprio, disponível no
repositório~\cite{mqttsecure2026repo}, que remove os tipos de documento não
admitidos pelo CI2 (normas, revistas de divulgação e capítulos de livro) e os
duplicados, por DOI e, na falta deste, por título normalizado. A
Tabela~\ref{tab:resultados-bases} regista os registos devolvidos por cada base.

\begin{table}[ht]
  \centering
  \caption{Registos devolvidos pela \emph{string} principal em cada base de dados.}
  \label{tab:resultados-bases}
  \small
  \begin{tabular}{@{}lr@{}}
    \toprule
    \textbf{Base de dados} & \textbf{Registos} \\
    \midrule
    IEEE Xplore          & \nIEEE \\
    Scopus               & \nScopus \\
    ScienceDirect (duas consultas) & \nScienceDirect \\
    \midrule
    \textbf{Total}       & \textbf{\nTotalIdentificados} \\
    \bottomrule
  \end{tabular}
\end{table}

% ---------------------------------------------------------------------------
\subsection{Critérios de inclusão e exclusão}
\label{sec:criterios}

Um estudo é incluído se satisfizer \emph{todos} os critérios de inclusão e
\emph{nenhum} dos de exclusão (Tabela~\ref{tab:criterios}). A versão em conferência do
protocolo baseado em AugPAKE (\emph{Augmented Password-Authenticated Key
Exchange})~\cite{calabretta2018token} foi considerada elegível na
triagem, mas analisou-se apenas a versão em revista~\cite{calabretta2018mqttauth},
que os autores apresentam como alargada a partir dela~\cite[§I]{calabretta2018mqttauth}
(CE2).

\begin{table}[ht]
  \centering
  \caption{Critérios de inclusão (CI) e de exclusão (CE).}
  \label{tab:criterios}
  \small
  \begin{tabularx}{\textwidth}{@{}L{1cm}Y@{}}
    \toprule
    \textbf{ID} & \textbf{Critério} \\
    \midrule
    CI1 & Propõe, implementa ou avalia um mecanismo de segurança para MQTT ou
          MQTT-SN, ou analisa ameaças específicas do MQTT. \\
    CI2 & É um estudo primário com revisão por pares (revista, conferência ou
          \emph{workshop}). \\
    CI3 & Foi publicado entre 2014 e 2026. \\
    CI4 & O texto integral está acessível, em inglês ou português. \\
    CE1 & O MQTT é apenas usado como transporte, sem contributo de segurança. \\
    CE2 & É uma versão preliminar de outro estudo incluído. \\
    CE3 & É literatura cinzenta (tese, relatório técnico, patente). \\
    CE4 & É uma revisão (analisada à parte, Secção~\ref{sec:relacionados}). \\
    CE5 & Tem menos de quatro páginas ou não descreve o mecanismo com detalhe
          suficiente. \\
    CE6 & Trata só de outro protocolo (CoAP, AMQP, DDS). \\
    \bottomrule
  \end{tabularx}
\end{table}

% ---------------------------------------------------------------------------
\subsection{Seleção dos estudos e avaliação da qualidade}
\label{sec:selecao}

\paragraph{Triagem.} Removidos os duplicados e os tipos não admitidos, os
\nTriados{} registos foram classificados por título e resumo por um único revisor,
aplicando os critérios que não exigem o texto integral (CI1--CI3, CE1, CE3, CE4 e
CE6); em caso de dúvida, o registo seguia para a fase seguinte. Por limitação de
tempo, não houve um segundo revisor independente, pelo que não se calculou a
concordância entre revisores; esta limitação é discutida na
Secção~\ref{sec:discussao}. A triagem deixou \nElegiveis{} estudos elegíveis.

\paragraph{Seleção para análise em profundidade.} O número de estudos elegíveis é
incompatível com a leitura integral de todos no prazo desta primeira parte do
projeto. Adotou-se, por isso, um desvio ao protocolo, que se declara: analisaram-se em
texto integral \nTextoIntegral{} estudos identificados numa fase exploratória,
anterior às pesquisas formais. Todos foram devolvidos pelas pesquisas e
classificados como elegíveis na triagem, mas não foram escolhidos a partir da
lista de elegíveis segundo um critério definido \emph{a priori}. Um deles foi excluído por CI4, por o texto integral não estar acessível,
pelo que o conjunto analisado tem \nIncluidos{} estudos. Esta escolha é uma ameaça à
validade da síntese, discutida na Secção~\ref{sec:discussao}, e a lista completa
das decisões de triagem está no repositório.

\paragraph{Leitura integral.} Os estudos analisados foram lidos em texto integral,
orientando a leitura pelos itens de extração da Secção~\ref{sec:extracao}; a
secção e a página que sustentam cada classificação ficaram registadas no ficheiro
de rastreabilidade do repositório.

\paragraph{\emph{Snowballing}.} A revisão de Hintaw et al.~\cite{hintaw2023mqtt},
citada na comparação de revisões de Lakshminarayana et
al.~\cite[§I]{lakshminarayana2024securing}, não foi devolvida pelas bases e foi
acrescentada às revisões relacionadas por \emph{snowballing} retrospetivo.
A Figura~\ref{fig:prisma} resume o processo.

\begin{figure}[ht]
  \centering
  \begin{tikzpicture}[
      caixa/.style={draw=iscte, fill=fundotab, rounded corners=2pt, text width=7.4cm,
                    minimum height=1.0cm, align=center, font=\small},
      lado/.style={draw=cinza, fill=white, rounded corners=2pt, text width=5.2cm,
                   minimum height=1.0cm, align=center, font=\small},
      fase/.style={draw=iscte, fill=iscte, text=white, font=\tiny\bfseries,
                   rotate=90, minimum width=1.5cm, minimum height=0.5cm},
      seta/.style={-{Stealth[length=2mm]}, thick, iscte}]
    \node[caixa] (id) {Registos identificados nas bases de dados\\
      ($n = \nTotalIdentificados$)\\[-1pt]
      {\scriptsize IEEE \nIEEE; Scopus \nScopus; ScienceDirect \nScienceDirect}};
    \node[lado, right=0.6cm of id] (dup) {Removidos antes da triagem:\\
      duplicados ($n = \nDuplicados$)\\tipo não admitido ($n = \nRemovidosTipo$)};
    \node[caixa, below=0.9cm of id] (tri) {Registos triados por título e resumo\\($n = \nTriados$)};
    \node[lado, right=0.6cm of tri] (exta) {Excluídos (CI/CE)\\($n = \nExcluidosTA$)};
    \node[caixa, below=0.9cm of tri] (ti) {Estudos elegíveis ($n = \nElegiveis$)\\
      analisados em texto integral, da fase exploratória ($n = \nTextoIntegral$)};
    \node[lado, right=0.6cm of ti] (exti) {Excluídos no texto integral: CI4\\($n = \nExcluidosTI$)};
    \node[caixa, below=0.9cm of ti, fill=iscte!15] (inc) {\textbf{Estudos primários analisados}\\($n = \nIncluidos$)};
    \node[lado, right=0.6cm of inc] (sec) {Revisões relacionadas: \nSecundarios{}\\
      (\nSecundariaUnicos{} devolvidas pela \emph{string} secundária; 1 por \emph{snowballing})};
    \draw[seta] (id) -- (tri);
    \draw[seta] (tri) -- (ti);
    \draw[seta] (ti) -- (inc);
    \draw[seta] (id) -- (dup);
    \draw[seta] (tri) -- (exta);
    \draw[seta] (ti) -- (exti);
    \node[fase, left=0.25cm of id] {Identificação};
    \node[fase, left=0.25cm of tri] {Triagem};
    \node[fase, left=0.25cm of ti] {Elegibilidade};
    \node[fase, left=0.25cm of inc] {Inclusão};
  \end{tikzpicture}
  \caption{Diagrama de fluxo PRISMA 2020 da seleção dos estudos.}
  \label{fig:prisma}
\end{figure}

% ---------------------------------------------------------------------------
\subsection{Extração de dados}
\label{sec:extracao}

Para cada estudo analisado registaram-se os seguintes itens: ano e tipo de publicação
(QE1, QE2); tipo de contributo e categoria da taxonomia (QF1, QF2, QE3); ponto de
atuação e primitivas usadas (QF2); propriedades garantidas e confiança no
\emph{broker} (QF3, QE3); tipo de avaliação e uso de dispositivos restritos (QF4);
e limitações declaradas (QF5). O resumo da extração consta do
Anexo~\ref{anx:extracao}.
